Thus it suffices to prove that
\(\mathrm{pr}_2((a\times U')\cap V)\)
is dense in \(G\), that is, nonempty,
which is clear since \(a\in U'=\mathrm{pr}_1 V\).
\end{proof}

\sgasection{3}{Construction of a group from a rational law}
\numpara{3.0}\label{XVIII.3.0}
Suppose that a prescheme \(X/S\) and a rational map
\(X\times_S X\to X\) over \(S\) are given,
and seek a group \(G/S\) and a map birational relative
to \(S\)\nde{Make explicit the notion of birational
map (relative to \(S\)), perhaps in an addition 1.5.2?}
\[
X\dashrightarrow G
\]
which commutes with the composition laws.
We treat only the case where \(X/S\) satisfies
the following hypothesis:
\[
\begin{gathered}
X/S\text{ is faithfully flat, of finite presentation,}\\
\text{and has separated fibers \og without embedded components\fg.}
\end{gathered}\tag{$\diamond$}
\]
(Note that the last two hypotheses are properties
true for a group prescheme.\nde{Cf. VIA,
to be specified\ldots})

We shall often suppress the symbol \(S\)
in fibered products.
Let \(X/S\) be a prescheme having properties
(\(\diamond\)) above, and let \(W\) be a
subprescheme of finite presentation of
\(X\times X\times X\) having the following property:
\[
\begin{gathered}
\text{The three morphisms }W\to X\times X
\text{ given by the projections of }X^3\text{ onto }X^2\\
\text{are open immersions schematically dense relative to }S.
\end{gathered}\tag{*}
\]
\begin{enoncerm}{Notations}{}
We shall use the following terminology:
given sections \(a,b,c\in X(S')\), with
\(S'\to S\) arbitrary, such that
\((a,b,c)\in W(S')\), we write:
\[
c=ab,\qquad b=a^{-1}c,\qquad a=cb^{-1}.
\]
\end{enoncerm}
\begin{definition}{3.0.1}\label{XVIII.3.0.1}
\nde{The numbering 3.0.1 and 3.0.2 has been introduced.}
We say, given a section \((a,b)\in X^2(S')\),
that \(ab\) is defined if and only if
there exists a section \(c\in X(S')\)
such that \((a,b,c)\in W(S')\),
i.e. if and only if \((a,b)\) is in
\(\mathrm{pr}_{12}W(S')\).
Similarly, saying that \(a^{-1}b\) or
\(ab^{-1}\) is defined has the analogous
meaning, and one also extends this
terminology to products of several factors.
% typo: "a la significations" -> "a la signification".
\end{definition}
\begin{remark}{3.0.2}\label{XVIII.3.0.2}
Let us immediately note the following fact:
by (i), \(W\) defines a rational map
\(X^2\to X\) over \(S\) (that given by
\((a,b)\mapsto ab\)).
It may well happen that this rational
map has a domain of definition larger
than \(\mathrm{pr}_{12}W\).
Nevertheless, we say that \(ab\) is
defined only if \((a,b)\in\mathrm{pr}_{12}W(S')\).
% typo?: reference (i) retained as printed.
\end{remark}
\begin{definition}{3.1}\label{XVIII.3.1}
A \emph{group germ} is a prescheme \(X/S\)
having properties (\(\diamond\)) above,
and a subprescheme \(W\) of \(X\times X\times X\)
of finite presentation having the following properties:

(i) \(W\) satisfies property (*) above.

(ii) For every section \(a\in X(S)\), the sets
\[
\begin{aligned}
\mathrm{pr}_i((a\times X\times X)\cap W),&\qquad i=2,3,\\
\mathrm{pr}_i((X\times a\times X)\cap W),&\qquad i=3,1,\\
\mathrm{pr}_i((X\times X\times a)\cap W),&\qquad i=1,2,
\end{aligned}
\]
are schematically dense in \(X\)
relative to \(S\), and this statement
remains true after every base change
\(S'\to S\).
(Hypothesis (i) implies that these
sets are open subsets of \(X\).)
Intuitively, this means that
\og \(ax\) is defined for sufficiently
general \(x\)\fg, etc.

(iii) The law is \emph{associative},
that is, if \((a,b,c)\in X^3(S')\),
with \(S'\to S\) arbitrary, is such
that \((ab)c\) and \(a(bc)\) are
defined, one has \((ab)c=a(bc)\).
\end{definition}
\begin{remarks}{}
(a) In (i), one may replace the
condition that \(W\) be schematically
dense in \(X^2\) relative to \(S\)
by the condition that, for every
\(s\in S\), the fiber \(W_s\) be
dense in \(X_s^2\) in the topological
sense, thanks to the hypotheses on \(X\) and to 1.2.

(b) Condition (iii) is equivalent
to the following:
let \(V_1\) (resp. \(V_2\)) be the
open subset of \(X^3\) where the
rational map \((a,b,c)\mapsto a(bc)\)
(resp. \((a,b,c)\mapsto(ab)c\))
is defined.
{\renewcommand{\thefootnote}{(*)}\addtocounter{footnote}{-1}%
\footnote{In the sense explained in 3.0.1;
one could also replace these open subsets
by the domains of definition (cf. 1.5)
of the rational maps under consideration.}}
Then there exists an open subset
\(V\subseteq V_1\cap V_2\) schematically
dense in \(X^3\) relative to \(S\)
such that the two preceding maps
agree on \(V\).
This is a consequence of 1.4,
since, of course, \(V_1,V_2\)
are schematically dense in \(X^3\)
relative to \(S\).

(c) Hypothesis (ii) will serve
below to ensure that \(X\) will
be a subobject of the group
prescheme which it defines.
In many cases, one can deduce
(ii) from (i), provided that
one replaces \(X\) by a relatively
schematically dense open subset
\(X'\), and \(W\) by a relatively
schematically dense open subset
of \(W\cap X'^3\).
In fact, one has:
\end{remarks}
\begin{proposition}{3.2}\label{XVIII.3.2}
Suppose that each geometric fiber
of \(X/S\) is irreducible, and
let \(W\) be a subprescheme of
\(X^3\) satisfying condition
(i) of 3.1.
Then there exists a relatively
schematically dense open subset
\(X'\) of \(X\), and a relatively
schematically dense open subset
\(W'\) of \(W\cap(X'\times X')\),
such that the pair \((X',W')\)
satisfies conditions (i) and (ii).
If (iii) holds for \((X,W)\),
it holds for \((X',W')\).
% typo?: W intersect (X' times X') retained as printed.
\end{proposition}
\begin{proof}
Set \(X'=\bigcap_{i=1}^3\mathrm{pr}_i W\).
Each \(\mathrm{pr}_i W\) is open
in \(X\), since \(W\to X^2\)
is an open immersion, and the
projections \(X^2\to X\) are
flat and of finite presentation.
Moreover, \(\mathrm{pr}_i W\)
is relatively schematically
dense in \(X\), since \(W\)
is so in \(X^2\), and the
projections \(X^2\to X\) are
surjective (it suffices to
verify density in the topological
sense). Take \(W'=W\cap X'^3\).

To verify that (i) is true, note that
\[
\begin{aligned}
W'&=(W\cap(X'\times X'\times X))
\cap(W\cap(X\times X'\times X'))\\
&\quad\cap(W\cap(X'\times X\times X')).
\end{aligned}
\]
Now \(W\cap(X'\times X'\times X)
\simeq\mathrm{pr}_{12}W\cap(X'\times X')\),
hence is relatively schematically
dense in \(W\).
Similarly, the other terms of
the right-hand member are
relatively schematically dense
in \(W\), and consequently
\(W'\) is relatively schematically
dense in \(W\).
Thus \(\mathrm{pr}_{ij}W'\)
is relatively schematically
dense in \(\mathrm{pr}_{ij}W\),
hence in \(X\times X\), hence
in \(X'\times X'\).

To verify condition (ii), let
\(a\in X'(S')\), and write \(S'=S\).
We must prove that, for example,
\(\mathrm{pr}_2((a\times X'\times X')\cap W')\)
is schematically dense in
\(X'\) relative to \(S\).
By 1.3, it suffices to verify
this fiber by fiber, that is,
it suffices to treat the case
where \(S\) is the spectrum
of a field, and in this case
it suffices to verify that
the open subset is nonempty,
since the fibers of \(X/S\)
are irreducible and \og without
embedded components\fg.
Since \(\mathrm{pr}_2((a\times X\times X)\cap W)\)
is nonempty, \(a\) being a
section of \(X'\), this open
subset is dense in \(X\).
One has
\[
\mathrm{pr}_2((a\times X\times X)\cap W)
=\mathrm{pr}_2((a\times X)\cap\mathrm{pr}_{12}W).
\]
Thus
\[
\begin{aligned}
\mathrm{pr}_2((a\times X)\cap\mathrm{pr}_{12}W)\cap X'
 &=\mathrm{pr}_2((a\times X')\cap\mathrm{pr}_{12}W)\\
 &=\mathrm{pr}_2((a\times X'\times X)\cap W)
\end{aligned}
\]
is dense in \(X\), hence in \(X'\).
Similarly, \(\mathrm{pr}_3(a\times X\times X'\cap W)\)
is dense in \(X\), hence in
\(\mathrm{pr}_3(a\times X\times X\cap W)\);
that is, \((a\times X\times X'\cap W)\)
is dense in \((a\times X\times X\cap W)\);
that is, \(\mathrm{pr}_2(a\times X\times X'\cap W)\)
is dense in \(\mathrm{pr}_2(a\times X\times X\cap W)\),
hence dense in \(X\), hence dense in \(X'\).
Now, since
\[
\begin{aligned}
\mathrm{pr}_2(a\times X'\times X'\cap W')
 &=\mathrm{pr}_2(a\times X'\times X'\cap W)\\
 &=\mathrm{pr}_2(a\times X'\times X\cap W)
\cap\mathrm{pr}_2(a\times X\times X'\cap W),
\end{aligned}
\]
it is indeed dense in \(X'\),
which is what had to be proved.
The other assertions of (ii)
follow by symmetry, and the
fact that condition (iii) is
preserved is trivial.
\nde{The following sentence has been added.}
Proposition 3.2 is proved.
\end{proof}

\nde{The numbering 3.2.1, as well as 3.2.2, has been added.}
\numpara{3.2.1}\label{XVIII.3.2.1}
Now fix a group germ \((X,W)\)
over \(S\).
We must make some preliminary
remarks on the situation, which
we have gathered below.
We shall often use these rules
without explicit mention in what follows.

Let \(a\in X(S')\), and write \(S'=S\).
Then one obtains a (bi)rational
map \(\varphi\) over \(S\) from
\(X\) to itself by associating
to a section \(x\) the section
\(ax\), if it is defined.
By 3.1 (ii), the domain of
definition of \(\varphi\)
contains the relatively
schematically dense open subset
\(\mathrm{pr}_2((a\times X\times X)\cap W)\),
and \(\varphi\) defines an
isomorphism from this open
subset onto the open subset
(where \(a^{-1}x\) is defined)
\(\mathrm{pr}_3((a\times X\times X)\cap W)\).
This remark is generalized
in an evident manner in Rule 1.

\begin{enoncerm}{Rule}{1}\label{XVIII.rule.1}
Let \(P=P(x,t_1,\ldots,t_n)\)
be a \og product\fg{} of the
symbols \(x,t_1,\ldots,t_n\),
obtained recursively as follows:
\(P_0=x\); \(P_{i+1}\) is one
of the following expressions:
\[
P_it,\quad tP_i,\quad P_i^{-1}t,\quad
tP_i^{-1},\quad P_it^{-1},\quad t^{-1}P_i,
\]
where \(t\) is one of the \(t_j\);
\(P_r=P\).
Let \(a_1,\ldots,a_n\in X(S')\).
Then there exists a relatively
schematically dense open subset
\(U\) of \(X_{S'}\) such that
the product \(P(x,a_1,\ldots,a_n)\)
is defined (in the sense of
Remark 3.0.2) for a section
\(x\in X(S'')\) if and only
if \(x\in U(S'')\), and the
map \(x\mapsto P(x,(a))\)
gives an isomorphism from
\(U\) onto another relatively
schematically dense open subset
of \(X_{S'}\), denoted \(P(U,(a))\).
\end{enoncerm}
\begin{enoncerm}{Rule}{2}\label{XVIII.rule.2}
Let \(a,b\in X(S')\). Then:

If \(ab\) is defined, so is
\(a^{-1}(ab)\), and \(a^{-1}(ab)=b\).

If \(a^{-1}b\) is defined,
so is \(a(a^{-1}b)\), and \(a(a^{-1}b)=b\).

If \(ba^{-1}\) is defined,
so is \((ba^{-1})a\), and \((ba^{-1})a=b\).
\end{enoncerm}
\begin{enoncerm}{Rule}{3}\label{XVIII.rule.3}
Let \(a,b,b'\in X(S')\).
If \(ab=ab'\), if \(ba=b'a\),
if \(a^{-1}b=a^{-1}b'\), or
if \(ba^{-1}=b'a^{-1}\), then \(b=b'\).
Here it is understood that
the equality relation implies
that both sides are defined.
\end{enoncerm}
\begin{enoncerm}{Rule}{4}\label{XVIII.rule.4}
Let \(a,b,c\in X(S')\).
Then, whenever both sides are
defined, one has
\[
a((ba)^{-1}c)=b^{-1}c.
\]
Similarly:
\[
\begin{gathered}
(c(ab)^{-1})a=cb^{-1},\qquad
a^{-1}((ab^{-1})c)=b^{-1}c,\\
(c(b^{-1}a))a^{-1}=cb^{-1}.
\end{gathered}
\]
\end{enoncerm}
\begin{enoncerm}{Rule}{5}\label{XVIII.rule.5}
All the following associativity
laws are true whenever both
sides are defined:
\[
\begin{gathered}
(a^{-1}b)c=a^{-1}(bc),\qquad
(ab^{-1})c=a(b^{-1}c),\qquad
(ab)c^{-1}=a(bc^{-1}),\\
(ab)^{-1}c=b^{-1}(a^{-1}c),\qquad
(a^{-1}b)c^{-1}=a^{-1}(bc^{-1}),\\
(ab^{-1})c^{-1}=a(cb)^{-1}.
\end{gathered}
\]
\end{enoncerm}
\numpara{3.2.2}\label{XVIII.3.2.2}
\textbf{Verification of the rules.}
(1) is proved by evident induction
on the length \(r\) of \(P\),
the case \(r=1\) being a direct
consequence of (3.1) (ii).

(2) This is trivial by definition.

(3) Indeed, by Rule 2, for
example in the first case,
one has
\[
b=a^{-1}(ab)=a^{-1}(ab')=b'.
\]
% typo: "ona" -> "on a".

(4) Let us verify, for example,
the first relation:
on the right-hand side,
left multiplication by \(b\)
is defined and gives \(c\),
by Rule 2.
Suppose that it is also
defined on the left-hand side.
Then one will have
\[
b(a((ba)^{-1}c))=(ba)((ba)^{-1}c)=c.
\]
Indeed, \((ba)\) is defined
by hypothesis, since it occurs
in the expression.
Thus the middle member is
defined and equal to \(c\)
by Rule 2, and equal to
the left-hand member by
associativity (cf. 3.1 (iii)).
Thus Rule 3 implies that
the desired equality is true
if this multiplication by
\(b\) is defined.
Now fix \(a,b\).
Then Rule 1 implies that
\(b(a((ba)^{-1}c))\) is indeed
defined for \(c\) \og in\fg{}
an open subset \(U\) of \(X\)
relatively schematically dense;
hence, in this relatively
schematically dense open subset,
the two rational maps
\(c\mapsto b^{-1}c\) and
\(c\mapsto a((ba)^{-1}c)\)
are equal.
By 1.4, they are equal on
every common domain of
definition, whence the desired result.

(5) This is the same kind
of argument as the preceding one.
For example, one verifies
\((ab^{-1})c^{-1}=a(cb)^{-1}\)
as follows:
if right multiplication by
\(c\) is defined on the
right-hand side, one has
equality by Rule 4.
Since it suffices to verify
such a formula on a relatively
schematically dense open subset,
one is reduced to the case
where this multiplication is
indeed defined.

\numpara{3.2.3}\label{XVIII.3.2.3}
Now consider the relation
\(R\) on \(X\times X\)
obtained by setting, for
\(a,b,a',b'\in X(S')\),
\((a,b)\sim(a',b')\) modulo
\(R(S')\) if and only if
there exists an \(S'\to S\)
covering for fppf, and a
section \(x\in X(S')\) such
that \((xa)b,(xa')b'\) are
defined and equal.
Then \(R\) is an equivalence relation.

Indeed, this relation is
evidently symmetric.
By Rule 1, the product
\((xa)b\) is defined if
\(x\) is \og in\fg{} a
suitable relatively schematically
dense open subset.
Thus 1.7 asserts that
there exists an \(S'\to S\)
covering for fppf, and
an \(x\in X(S')\), such
that \((xa)b\) is defined.
The relation is therefore
reflexive, and transitivity
is a consequence of the
following lemma:
\begin{lemma}{3.3}\label{XVIII.3.3}
Let \(x,y,a,b,a',b'\) be
sections of \(X(S')\) such
that \((xa)b,(xa')b',(ya)b,(ya')b'\)
are defined.
If \((xa)b=(xa')b'\),
then \((ya)b=(ya')b'\).
\end{lemma}
Indeed, the lemma says
that one can test
\((a,b)\sim(a',b')\) with
an arbitrary \(x\in X(S')\),
\(S'\to S\) fppf-covering,
such that both products
are defined.
Given \(a,b,a',b',a'',b''\in X(S)\),
one may, by Rule 1, 1.1,
and 1.7, find an fppf-covering
extension \(S'\to S\) and
a section \(x\in X(S')\)
such that the three products
in question are defined,
whence transitivity.
\begin{proof}[Proof of the lemma]
To begin with, write formally:
\[
\begin{aligned}
(za)b
 &=(((zx^{-1})x)a)b\\
 &=((zx^{-1})(xa))b=(zx^{-1})((xa)b),\\
(za')b'
 &=(((zx^{-1})x)a')b'\\
 &=((zx^{-1})(xa'))b'=(zx^{-1})((xa')b').
\end{aligned}
\]
One verifies that these
equalities are indeed true
if the members are defined,
by the appropriate rules.
\nde{To be specified\ldots}
It follows that
\((za)b=(za')b'\) if
all these expressions are
defined.
Moreover, by Rule 1 and
the hypotheses already made,
these expressions are indeed
defined if \(z\in X(S'')\)
is in \(V(S'')\), where
\(V\) is a certain relatively
schematically dense open
subset of \(X\) (we have
taken \(S'=S\)).
Thus the two rational
maps from \(X\) to itself
given by \(z\mapsto(za)b\)
and \(z\mapsto(za')b'\)
are equal, whence
\((ya)b=(ya')b'\).
\end{proof}
\begin{lemma}{3.4}\label{XVIII.3.4}
Consider the rational map
\(\varphi:X^3\to X\)
over \(S\), defined by
\((a,b,c)\mapsto c^{-1}(ab)\).
Let \(U\) be the domain
of definition of \(\varphi\),
and consider the graph
\(\Gamma\) of the morphism
\(f:U\to X\) induced by
\(\varphi\), which is a
subscheme of \(X^4\).
Then a section
\((a,b,c,d)\in X^4(S)\)
is in \(\Gamma(S)\)
if and only if \((a,b)\sim(c,d)\).
\end{lemma}
\begin{proof}
First note that the
rational map \(\varphi\)
is the same as that given
by the formula
\((a,b,c)\mapsto(xc)^{-1}((xa)b)\)
for an arbitrary section
\(x\in X(S)\).
It amounts to the same
to say that
\(c^{-1}(ab)=(xc)^{-1}((xa)b)\)
whenever both sides are defined.
We leave the verification
to the reader.

Let us then show that
the map \(\varphi\) is
defined for a section
\((a,b,c)\in X^3(S)\)
if and only if there
exists \(d\) with
\((a,b)\simeq(c,d)\),
and that then
\(d=\varphi(a,b,c)\).
Indeed, suppose that
one has \((a,b)\sim(c,d)\).
To verify that the map
\(\varphi\) is defined,
one may make an fppf-covering
base extension (Proposition 1.6),
and may therefore suppose
that there exists a
section \(x\in X(S')\)
such that \((xa)b,(xc)d\)
are defined and equal.
It follows (Rule 2) that
\((xc)^{-1}((xa)b)\)
is defined and equal to \(d\).

Conversely, suppose that
the map \(\varphi\) is
defined at the section
\((a,b,c)\), and let
\(d=\varphi(a,b,c)\).
Choose \(S'\to S\)
fppf-covering and a
section \(x\in X(S')\)
such that \((xa)b,(xc)d\)
are defined.
We wish to show that
they are equal.
For this, it suffices
to prove that the
two rational maps over
\(S\) from \(X\) to
itself, given by
\(b\mapsto(xa)b\)
and \(b\mapsto(xc)\varphi(a,b,c)\),
are the same, which
follows from the remark
in the first paragraph.
% typo: "définieen" -> "définie en".
\end{proof}
\begin{proposition}{3.5}\label{XVIII.3.5}
The equivalence relation
\(R\rightrightarrows X^2\)
is representable, and
is a flat relation of
finite presentation, that
is, the projections
of \(R\) onto \(X^2\)
are flat morphisms of
finite presentation.
\end{proposition}
\begin{proof}
One may suppose \(S\) affine.
Since \((X,W)\) is of
finite presentation over
\(S\), one can descend
the entire situation to
an \(S\) of finite type
over \(\Spec\ZZ\),
hence noetherian.
We may therefore suppose
that \(S\) is noetherian.
Then it is trivial that
the graph \(\Gamma\)
of (3.4) is of finite
presentation over \(X^2\).
The projection
\(\mathrm{pr}_{12}|_\Gamma\)
is flat, since
\(\Gamma\isomto\mathrm{pr}_{123}\Gamma\),
which is an open subset
of \(X^3\), and the
projection
\(\mathrm{pr}_{12}:X^3\to X^2\)
is flat because \(X\)
is flat over \(S\).

I say that \(\Gamma\)
represents \(R\).
Note that there is
something to prove,
since the domain of
definition of a rational
map does not in general
commute with base extensions.
Let \(S''\to S\).
What is clear, by
(3.4) applied to \(S''\),
is that
\(R(S'')\supset\Gamma(S'')\),
since \(\varphi\times_S S''\)
is certainly defined on
\(U\times_S S''\).
Thus let
\((a,b,c,d)\in R(S'')\).
One must prove that
\((a,b,c,d)\in\Gamma(S'')\).
The verification is
made locally for, say,
the étale topology.
One may therefore suppose
that \(S''\) is strictly
local, i.e. the spectrum
of a henselian ring
with separably closed
residue field.
Moreover, applying (1.6)
and the customary sorites
of passage to the limit,
one is reduced to the
case where \(S\) is
strictly local and
\(S''\to S\) is local.
Suppose that one has
a section \(x\in X(S)\)
such that over \(S''\)
the products \((xa)b,(xc)d\)
are defined.
This will imply that
\((xc)^{-1}((xa)b)\)
is defined and equal to \(d\).
Now there exists a
relatively schematically
dense open subset \(V\)
of \(X^3\) such that
\((xc)^{-1}((xa)b)\)
is defined if and
only if
\((a,b,c)\in V(S'')\),
and one has \(V\subseteq U\).
Thus
\((a,b,c,d)\in\Gamma(S'')\)
if such an \(x\) exists.
By (1.6), one may
make an fppf-covering
base extension \(S'\to S\)
to find such an \(x\).
Since \(S\) is strictly
local, one can
{\renewcommand{\thefootnote}{(*)}\addtocounter{footnote}{-1}%
\footnote{Combining Exp. VIA, 1.1.1
and EGA IV$_4$, 17.16.2.}}
find \(S'\to S\)
faithfully flat, local,
and finite, and a
section \(x\in X(S')\)
which \og passes through\fg{}
an arbitrary closed
point of the closed
fiber of \(X/S\).
